\noindent\textit{2015 manuscript.}

\section{Thermodynamic potentials}
\noindent\textit{In the 2015 manuscript these subsections stood in the section Reactions, after the Gibbs sum.}

\subsection{Gibbs Free Energy}

cf. Section ``Gibbs Free Energy'' of Ch. 9 Gibbs Free Energy and Chemical Reactions from Kittel and Kroemer (1980) \cite{CKittelHKroemer1980}

Recall 
\[
\begin{aligned}
  & F = U- \tau \sigma \\ 
  & dF = dU - \tau d\sigma - \sigma d\tau = - \sigma d\tau -pdV + \mu dN
\end{aligned}
\]
where the \emph{generalized} thermodynamic identity
\[
dU = Q + W + \mu dN = \tau d\sigma - pdV + \mu dN
\]
was used. 

Now
\[
\begin{aligned}
  & G \equiv U - \tau \sigma + pV = F + pV \\ 
  & dG = dU - \tau d\sigma - \sigma d\tau + pdV + Vdp = \mu dN - \sigma d\tau + Vdp
\end{aligned}
\]
and so $G = G(N,\tau , p)$.  

Now \[
\begin{gathered}
  \sigma = \frac{1}{\tau} [ U +pV - G ] \\
  \Longrightarrow d\sigma = \frac{-1}{\tau^2} d\tau ( U + pV - G) + \frac{1}{\tau} (dU + V dp + pdV - dG)
\end{gathered}
\]
Consider a curve $\gamma : \mathbb{R} \to \Sigma$ s.t. $dU(\dot{\gamma}) = dp(\dot{\gamma}) = dV(\dot{\gamma})=0$ corresponding to minimized internal energy, constant pressure, and constant volume.  Then
\[
\begin{gathered}
  d\sigma(\dot{\gamma}) = 0 + 0 + 0 + -dG(\dot{\gamma}) \geq 0 
\end{gathered}
\]
and so $dG(\dot{\gamma}) \leq 0$, so $G$ must be a minimum, because $d\sigma(\dot{\gamma}) \geq 0$ represents the fact that entropy must always increase for any thermodynamic process.  $G$ must move to a minimum at equilibrium, i.e. ``For any irreversible change taking place entirely within $\mathcal{S}$ will increase $\sigma$ and thus decrease $G_{\mathcal{S}}$'' \cite{CKittelHKroemer1980} for $\mathcal{S}$ signifying the system.  

$\tau,p$ are \textbf{intensive quantities}; they do not change value when 2 identical systems are put together.  

$U,\sigma,V,G$ are linear in $N$; their values doubles when 2 identical systems are put together; apart from interface effects \\
$U,\sigma,V,N,G$ are \textbf{extensive quantities}.  


\[
\begin{gathered}
G = N\varphi(\tau,p) \\
\left( \frac{ \partial G}{ \partial N} \right)_{p,\tau} = \varphi(\tau,p) = \mu \Longrightarrow G(N,\tau,p) = N\mu(\tau,p)
\end{gathered}
\]

If more than 1 chemical species present, 
\begin{equation}\label{Eq:Gibbspotential}
  \boxed{ 
    \begin{gathered}
      G = \sum_j N_j \mu_j  \\
      \tau d\sigma = dU + pdV - \sum_j \mu_j dN_j \\
      dG = \sum_j \mu_j dN_j - \sigma d\tau + Vdp 
      \end{gathered}
    }
\end{equation}


\subsection{Summary of Thermodynamic Potentials}

With Eq. \ref{Eq:Gibbspotential}, I will summarize the thermodynamic potentials we have, obtained through Legendre transformations:
\[
\begin{gathered}
  Q = dU - W = dU + p dV
\end{gathered}
\]
For curve or thermodynamic process in $\Sigma$, $\gamma$, s.t. $dp(\dot{\gamma})=0$ (thermodynamic process at \emph{constant pressure}),
\[
\begin{gathered}
  Q(\dot{\gamma}) = dH(\dot{\gamma}) -Vdp(\dot{\gamma}) = dH(\dot{\gamma}) - 0 \\
  \Longrightarrow \int_{\gamma} Q(\dot{\gamma}) = \int dH(\dot{\gamma}) = H_f - H_i \qquad \, \text{ for constant $p$ }
\end{gathered}
\]
with 
\[
\begin{gathered}
  H = U + pV \\
  dH = dU + pdV + Vdp = Q + Vdp
\end{gathered}
\]
For the thermodynamic potential
\begin{equation}\label{Eq:ThermodynamicPotentialsSummary}
\begin{gathered}
  \begin{gathered}
    dU = \tau d\sigma - pdV \\ 
    \boxed{ U = U(\sigma, V) } \\
    dU = \tau d\sigma - pdV + \sum_j \mu_j dN_j \qquad \, U = U(\sigma, V,N_j)
\end{gathered} \\
  \begin{gathered}
    F = U -\tau \sigma  \\
    dF = dU- \tau d\sigma - \sigma d\tau = -pdV - \sigma d\tau  \Longrightarrow \boxed{ F = F(\tau,V) } \\  
    dF = Q + W - d(\tau \sigma ) 
\end{gathered} \\
  \begin{gathered}
    G = F + pV \\ 
    dG = -\sigma d\tau + Vdp \Longrightarrow \boxed{ G = G(\tau,p) } \\
    dG = (Q - pdV ) + pdV + Vdp - d(\tau \sigma ) = Q + Vdp - d(\tau \sigma)
\end{gathered}
\end{gathered}
\end{equation}
Let $\Delta H := H_f - H_i$.  \\
If $\Delta H < 0$, $\int_{\gamma} Q(\dot{\gamma}) <0$, heat transferred to surroundings: reaction is \emph{exothermic}.  \\
If $\Delta H > 0$, $\int_{\gamma} Q(\dot{\gamma}) >0$, heat transferred to system: reaction is \emph{endothermic}.  
Now $dH = Q + Vdp$.  Suppose $Q = C_p(\tau) d\tau$.  then for constant presure, $dp(\dot{\gamma})=0$ and so 
\[
\int_{\gamma} dH = \Delta H = \int_{\gamma} C_p(\tau) d\tau
\]

For constant volume, $dV(\dot{\gamma}_V)=0$, and so 
\[
Q(\dot{\gamma}_V) \equiv Q_V = dU(\dot{\gamma}_V)
\]

\subsubsection{Relations between derivatives of thermodynamic quantities}

For an adiabatic process, $Q=0$. 

For $Q= Q(\tau, V)$, 
\begin{equation}\label{Eq:QAsA1Form}
Q = \tau d\sigma = C_V d\tau + \left( \frac{ \partial Q}{ \partial V} \right)_{\tau} dV
\end{equation}
and $\sigma = \sigma(\tau, V)$, so 
\begin{equation}\label{Eq:EntropyAsA1Form}
\tau d\sigma = \tau \left( \frac{ \partial \sigma}{ \partial \tau } \right)_V d\tau + \tau \left( \frac{ \partial \sigma}{ \partial V} \right)_{\tau} dV
\end{equation}
So by comparing terms between Eq. \ref{Eq:QAsA1Form} and Eq. \ref{Eq:EntropyAsA1Form}, conclude that 
\begin{equation}\label{Eq:CompareQAndEntropyTermsInTheDifferentialForms}
\left( \frac{ \partial Q}{ \partial V} \right)_{\tau} = \tau \left( \frac{ \partial \sigma}{\partial V} \right)_{\tau}
\end{equation} 
For curve $c= c(t) = (\tau(t), V(t))$ s.t. $\sigma = \sigma(\tau, V) = $ constant, 
\[
\begin{gathered}
 \dot{c} = \dot{\tau} \frac{ \partial }{ \partial \tau} + \dot{V} \frac{ \partial }{\partial V}  \\
 Q(\dot{c}) = 0 = C_V\dot{\tau} + \dot{V} \left( \frac{ \partial Q}{\partial V} \right)_{\tau} \text{ or } \dot{\tau} = \frac{- \dot{V}}{C_V} \left( \frac{\partial Q}{ \partial V} \right)_{\tau} 
\end{gathered}
\]

\quad \\ 
Let $t=V$. Let $\dot{\tau} = \left( \frac{\partial \tau}{\partial V} \right)_{\sigma}$, \, $\dot{V} = 1$.  So using Eq. \ref{Eq:CompareQAndEntropyTermsInTheDifferentialForms}, 
\begin{equation}\label{Eq:ConstantVolumeHeatCapacityDifferentialRelationship}
 \left( \frac{ \partial \tau}{ \partial V} \right)_{\sigma} = \frac{-\tau}{C_V} \left( \frac{ \partial \sigma}{ \partial V} \right)_{\tau}
\end{equation}

Likewise, for $Q= Q(\tau, p)$, $\sigma = \sigma(\tau, p)$, 
\[
Q = Q(\tau, p) = C_p d\tau + \left( \frac{\partial Q}{\partial p } \right)_{\tau} dp = \tau d\sigma = \tau \left( \frac{\partial \sigma}{\partial \tau} \right)_p d\tau + \tau \left( \frac{\partial \sigma}{ \partial p} \right)_{\tau} dp
\]
and repeating all steps for $p$, but instead using $V$ (what I call "label symmetry"), 
\begin{equation}\label{Eq:ConstantPressureHeatCapacityDifferentialRelationship}
\left( \frac{\partial \tau}{\partial p } \right)_{\sigma} = \frac{-\tau}{C_p} \left( \frac{\partial \sigma}{\partial p} \right)_{\tau}
\end{equation}

Divide Eq. \ref{Eq:ConstantPressureHeatCapacityDifferentialRelationship} by Eq. \ref{Eq:ConstantVolumeHeatCapacityDifferentialRelationship}:
\begin{equation}\label{Eq:ConstantPressureHeatCapacityRelationshipOverConstantVolume}
\frac{ \left( \frac{\partial \tau}{\partial p}\right)_{\sigma} }{ \left( \frac{\partial \tau}{\partial V}\right)_{\sigma} } = \frac{C_V}{C_p} \frac{ \left( \frac{\partial \sigma}{\partial p}\right)_{\tau} }{ \left( \frac{\partial \sigma}{\partial V}\right)_{\tau} }
\end{equation} 

Since $F= F(\tau, V)$, 
\[
\begin{gathered}
dF = -pdV - \sigma d\tau \\
\begin{aligned}
& \left(\frac{\partial F}{\partial \tau} \right)_V = -\sigma \\ 
& \left(\frac{\partial F}{\partial V} \right)_{\tau} = - p 
\end{aligned}
\end{gathered}
\]
and it is sufficient that $\left(\frac{\partial F}{\partial \tau} \right)_V$, $\left(\frac{\partial F}{\partial V} \right)_{\tau}$ are continuous on a neighborhood of $(\tau, V)$ and $\frac{\partial^2 F}{ \partial \tau \partial V}$ exists and is continuous on neighborhood of $(\tau, V)$, then
\[
\left(\frac{\partial^2 F}{\partial \tau \partial V} \right) = \left(\frac{\partial^2 F}{\partial V \partial \tau} \right)
\]
(cf. search for Clairaut's Thm., mean value Thm., or Fubini's Thm.)

\begin{equation}\label{Eq:2ndOrderPartialDerivativeHelmholtzRelationship}
\Longrightarrow \left( \frac{ \partial \sigma}{ \partial V} \right)_{\tau} = \left(\frac{\partial p}{\partial \tau } \right)_V \qquad \text{(from Helmholtz's free energy $F$)}
\end{equation}

\quad \\ 
Similarly, with the Gibbs free energy $G$ given in Eq. \ref{Eq:ThermodynamicPotentialsSummary}, \\
since $G=G(\tau, p)$, 
\[
dG = -\sigma d\tau + Vdp
\]
then 
\[
\begin{gathered}
	\begin{aligned}
		& \left( \frac{\partial G}{\partial \tau} \right)_p = -\sigma \\ 
		& \left( \frac{ \partial G}{\partial p} \right)_{\tau} = V
	\end{aligned} 
\end{gathered}
\]
so 
\begin{equation}\label{Eq:2ndOrderPartialDerivativeGibbsRelationship}
-\left( \frac{ \partial \sigma}{ \partial p} \right)_{\tau} = \left(\frac{\partial V}{\partial \tau } \right)_p \qquad \text{(from Gibbs free energy $G$)}
\end{equation}

Use Eqns. \ref{Eq:2ndOrderPartialDerivativeHelmholtzRelationship}, \ref{Eq:2ndOrderPartialDerivativeGibbsRelationship} in Eq. \ref{Eq:ConstantPressureHeatCapacityRelationshipOverConstantVolume}:

\begin{equation}\label{Eq:ConstantPressureHeatCapacityRelationshipOverConstantVolume2}
\frac{ \left( \frac{\partial \tau}{\partial p}\right)_{\sigma} }{ \left( \frac{\partial \tau}{\partial V}\right)_{\sigma} } = \frac{-C_V}{C_p} \frac{ \left( \frac{\partial V}{\partial \tau}\right)_p }{ \left( \frac{\partial p}{\partial \tau}\right)_V } 
\end{equation}

\quad \\ 
Consider an \emph{adiabatic} process, such as adiabatic compression. Then $Q = \tau d\sigma = 0$. So choose curves in $\Sigma$ s.t. $\sigma$ is constant along the curve.

\quad \\
Let $\sigma(\tau, V), d\sigma = \left( \frac{\partial \sigma}{\partial \tau} \right)_V d\tau + \left( \frac{\partial \sigma}{\partial V} \right)_{\tau} dV$. \\
Let $c=c(t)= (\tau, p) = (\tau(t), p(t))$ s.t. $\sigma$ constant. Then
\begin{equation}\label{Eq:AdiabaticDifferentialRelationship1}
\begin{gathered}
d\sigma(\dot{c}) = 0 = \left( \frac{\partial \sigma}{\partial \tau} \right)_V \dot{\tau} + \left( \frac{\partial \sigma}{\partial V} \right)_{\tau} \dot{p} \left( \frac{\partial V}{ \partial p } \right)_{\sigma} \Longrightarrow \\
\left( \frac{\partial V}{\partial p} \right)_{\sigma} = \frac{ - \dot{\tau} \left( \frac{\partial \sigma}{\partial \tau} \right)_V }{ \dot{p} \left( \frac{\partial \sigma}{\partial V} \right)_{\tau} }\xrightarrow{t = p} \frac{ - \left( \frac{\partial \tau}{\partial p} \right)_{\sigma} \left( \frac{\partial \sigma}{\partial \tau}\right)_V }{ \left( \frac{\partial \sigma}{\partial V}\right)_{\tau} } 
\end{gathered}
\end{equation}
 
\quad \\
Now let $c=c(t) = (\tau, V) = (\tau(t), V(t))$ s.t. $\sigma$ constant.

\[
d\sigma(\dot{c}) = 0 = \left( \frac{\partial \sigma}{\partial \tau} \right)_V \dot{\tau} + \left( \frac{\partial \sigma}{\partial V} \right)_{\tau} \dot{V} \text{ or } \frac{ \left( \frac{\partial \sigma}{\partial \tau}\right)_V }{ \left( \frac{\partial \sigma}{\partial V}\right)_{\tau} } = \frac{  - \dot{V}}{\dot{\tau}} \xrightarrow{ t = V} \frac{-1}{ \left( \frac{\partial \tau}{\partial V} \right)_{\sigma} }
\]
Then plugging in the previous expression into Eq. \ref{Eq:AdiabaticDifferentialRelationship1}, we obtain
\begin{equation}\label{Eq:AdiabaticCompressionOfV}
\Longrightarrow \left( \frac{\partial V}{ \partial p} \right)_{\sigma} = \frac{ \left( \frac{\partial \tau }{\partial p} \right)_{\sigma} }{ \left( \frac{\partial \tau}{\partial V} \right)_{\sigma} }
\end{equation}

Using label symmetry for $p, V \mapsto V, p$, then
\[
\begin{gathered}
\left( \frac{\partial p}{\partial V} \right)_{\sigma} =  \frac{ - \left( \frac{\partial \tau}{\partial V} \right)_{\sigma} \left( \frac{\partial \sigma}{\partial \tau}\right)_p }{ \left( \frac{\partial \sigma}{\partial p}\right)_{\tau} } = \frac{ \left( \frac{\partial \tau }{\partial V} \right)_{\sigma} }{ \left( \frac{\partial \tau}{\partial p} \right)_{\sigma} } \text{ since } \\
\frac{ \left( \frac{\partial \sigma}{\partial \tau}\right)_p }{ \left( \frac{\partial \sigma}{\partial p}\right)_{\tau} } = \frac{-1}{ \left( \frac{\partial \tau}{\partial p} \right)_{\sigma} }
\end{gathered}
\]

The previous expression and its derivation from first principles (namely that $d\sigma$ is a 1-form on 2-dim. $\Sigma$ and thermodynamic quantities are coordinates for $\Sigma$) leads me to strongly exhort the following: stop memorizing and using "hand-wavey" formulae that claim to transform partial derivatives "to different independent variables" such as the "Jacobians" in Landau and Lifshitz (1980) \cite{LLandauELifshitz1980}), pp. 53, Ch. 11 "Thermodynamic Quantities", Sec. 16 "Relations between the Derivatives of Thermodynmaic Quantities", Eq. (IV) and related:
\[
\frac{ \partial (u,v)}{ \partial (x,y)} = \frac{ \partial (u,v) }{ \partial (t,s)} \cdot \frac{\partial (t,s) }{ \partial (x,y) }
\]
or in Appendix 4, "Summary of Thermodynamic Relations", pp. 577 of Sabersky, et. al. (1998) \cite{SAHG1998}, which claim that the "rules of differential calculus" (???) allows for some to write, for $z = z(x,y)$, 
\[
\begin{gathered} 
dz = \left( \frac{ \partial z}{ \partial x} \right)_y dx + \left( \frac{ \partial z}{ \partial y} \right)_x dy \text{ or } \\
\left( \frac{ \partial x }{ \partial y} \right)_z \left( \frac{ \partial z }{ \partial x} \right)_y \left( \frac{ \partial y}{ \partial z }\right)_x = -1
\end{gathered}
\]
For the latter, the "differential", $dx$, $dy$ is ill-defined. It is \href{https://math.stackexchange.com/questions/2101896/what-is-the-meaning-of-a-differential-in-terms-of-an-exact-differential}{"an outdated concept from the time of Liebniz which was used to define derivatives and integrals before limits came along."}. For the former, the "Jacobian" is not general enough (at least for me), being only applicable to Euclidean spaces and is another unnecessary formula to memorize.

Instead, consider the coordinate transformation of 2-dim. manifold $\Sigma$:
\[
(\tau, V) \mapsto (\tau, p)
\]
Consider then that in general, $p = p(\tau, V)$, so that 
\[
dp = \left( \frac{ \partial p }{ \partial \tau} \right)_V d\tau + \left( \frac{ \partial p }{ \partial V} \right)_{\tau} dV
\]
Consider curve $c\subset \Sigma$, $c=c(t) = (\tau, p) = (\tau(t), p(t))$, so that 
\[
\dot{c} = ( \dot{\tau}, \dot{p}) = \dot{\tau} \frac{ \partial }{ \partial \tau } + \dot{p} \frac{ \partial }{ \partial p}
\]

Then
\begin{equation}\label{Eq:CoordinateTransformationOfPressurep}
dp(\dot{c}) = \dot{p} = \dot{\tau} \left( \frac{ \partial p }{ \partial \tau} \right)_V + \dot{p} \left( \frac{\partial p}{ \partial V}\right)_{\tau}  \left( \frac{ \partial V}{ \partial p } \right)_{\tau}
\end{equation}
where we used the fact that $dV \left( \frac{ \partial }{ \partial \tau } \right) = 0 $ ($\tau$ and $V$ are still independent coordinates for $\Sigma$) and \\
$dV\left( \frac{ \partial }{ \partial p} \right) = \left( \frac{ \partial V}{ \partial p } \right)_{\tau}$ since $\tau$ is held constant for the component of $\dot{c} \in \mathfrak{X}(\Sigma)$ in $p$ (i.e. $\frac{\partial }{ \partial p}$).

Consider letting parameter $t=V$ for $c$. Then
\[
\dot{\tau} = \left( \frac{ \partial \tau }{ \partial V } \right)_p, \qquad \dot{p} = \left( \frac{ \partial p }{ \partial V} \right)_{\tau}
\]

Consider letting $t=p$. Then $\dot{\tau} = \left( \frac{ \partial \tau}{ \partial p} \right)_V$ and $\dot{p} = 1$. $V$ is help constant for $\dot{\tau} = \left( \frac{ \partial \tau}{ \partial p} \right)_V$ because of how $p = p(\tau, V)$ for 
\[
dp = \left(\frac{ \partial p}{ \partial \tau} \right)_V d\tau + \left( \frac{ \partial p }{ \partial V} \right)_{\tau} dV
\]
so
\begin{equation}\label{Eq:CoordinateTransformationOfPressurep2}
\Longrightarrow 1 = \left( \frac{ \partial \tau }{ \partial p } \right)_V \left( \frac{ \partial p }{ \partial \tau } \right)_V + \left( \frac{ \partial p }{ \partial V } \right)_{\tau} \left( \frac{ \partial V }{ \partial p } \right)_{\tau}
\end{equation}

Plug in Eq. \ref{Eq:CoordinateTransformationOfPressurep2} into Eq. \ref{Eq:CoordinateTransformationOfPressurep}:
\[
\begin{gathered}
	\dot{p} = \dot{\tau} \left( \frac{ \partial p}{ \partial \tau} \right)_V + \dot{p} \left[ 1 - \left( \frac{ \partial \tau}{ \partial V} \right)_p \left( \frac{\partial p}{ \partial \tau} \right)_V \right] \Longrightarrow \dot{p} \left( \frac{ \partial \tau }{ \partial V } \right)_p \left( \frac{ \partial p }{ \partial \tau } \right)_V = \dot{\tau} \left( \frac{ \partial p }{ \partial \tau } \right)_V \text{ or }
\end{gathered}
\]
\begin{equation}
\boxed{ \left( \frac{ \partial p }{ \partial V} \right)_{\tau} = \frac{\left( \frac{ \partial \tau }{ \partial V } \right)_p }{ \left( \frac{ \partial \tau }{ \partial p } \right)_V} }
\end{equation}

Consider the coordinate transformation $(\tau, p) \mapsto (\tau, V)$. By "label symmetry" (i.e. substitute $p$ for $V$, $V$ for $p$, for example),
\begin{equation}
\boxed{ \left( \frac{\partial V}{ \partial p } \right)_{\tau} = \frac{ \left( \frac{ \partial \tau }{ \partial p } \right)_V}{ \left( \frac{ \partial \tau }{ \partial V } \right)_p } }
\end{equation}

In fact, furthermore, by label symmetry, consider $(\tau, p) \mapsto (\tau, \rho)$, $\rho \equiv $ density, then
\begin{equation}
\boxed{ \left( \frac{\partial p}{ \partial \rho } \right)_{\tau} = \frac{ \left( \frac{ \partial \tau }{ \partial \rho } \right)_p}{ \left( \frac{ \partial \tau }{ \partial p } \right)_{\rho} } }
\end{equation}

Compare the previous expressions to the equation before (16.14) and (16.14) of Landau and Lifshitz (1980) \cite{LLandauELifshitz1980}, pp. 54,
\[
\begin{gathered}
\left( \frac{\partial V}{ \partial  P } \right)_S = \frac{\partial (V, S) }{ \partial (P, S)} = \frac{\partial( V,S) / \partial (V, T) }{ \partial (P, S)/ \partial (P, T)} \cdot \frac{ \partial (V, T) }{ \partial (P ,T)} = \frac{(\partial S/ \partial T)_V }{ (\partial S / \partial T)_P} \cdot \left( \frac{ \partial V}{ \partial P} \right)_T \\
\left( \frac{\partial V}{ \partial  P } \right)_S = \frac{ C_V}{ C_p} \left( \frac{\partial V}{ \partial  P } \right)_T
\end{gathered}
\]

Similarly with Eq. (A-4.9a) on pp. 577 with Appendix 4 of  Sabersky, et. al. (1998) \cite{SAHG1998}:
\[
\left( \frac{ \partial p }{ \partial \rho} \right)_s = - \frac{ (\partial s / \partial \rho)_p }{ (\partial s / \partial p)_{\rho}}  = - \frac{ (\partial s / \partial T)_p (\partial T / \partial \rho)_p }{ (\partial s / \partial p)_{\rho} (\partial T / \partial p )_{\rho}} = \frac{ (\partial s / \partial T)_p }{ (\partial s / \partial T)_{\rho}} \left( \frac{ \partial p }{ \partial \rho} \right)_T
\]

The expression we'll need uses $V$:
\[
\begin{gathered}
V = V(\tau, p) \Longrightarrow \\
dV = \left( \frac{\partial V}{\partial \tau} \right)_p d\tau + \left( \frac{\partial V}{\partial p} \right)_{\tau} dp 
\end{gathered}
\]
Let $c=(\tau, p)$  be a curve for when $V$ is constant.

\[
\begin{gathered}
dV(\dot{c}) = 0 = \left( \frac{\partial V}{\partial \tau} \right)_p \dot{\tau} + \left( \frac{\partial V}{ \partial p } \right){\tau} \dot{p} \xrightarrow{ t = \tau} - \left( \frac{\partial V}{ \partial \tau} \right)_p = \left( \frac{\partial V}{\partial p} \right)_{\tau} \left( \frac{\partial p}{\partial \tau} \right)_V \text{ or } 
\end{gathered}
\]
\begin{equation}
\left( \frac{\partial V}{\partial p} \right)_{\tau} = -\frac{\left( \frac{\partial V}{\partial \tau} \right)_p }{ \left( \frac{\partial p}{\partial \tau} \right)_V}
\end{equation}

Using label symmetry for $V, p \mapsto p, V$, then also
\[
\left( \frac{\partial p}{\partial V} \right)_{\tau} = -\frac{\left( \frac{\partial p}{\partial \tau} \right)_V }{ \left( \frac{\partial V}{\partial \tau} \right)_p}
\]

Then using both Eq. \ref{Eq:AdiabaticCompressionOfV} and the previous equation in Eq. \ref{Eq:ConstantPressureHeatCapacityRelationshipOverConstantVolume2}, then we get 
\begin{equation}
\begin{gathered}
\boxed{ \left( \frac{\partial V}{\partial p }\right)_{\sigma} = \frac{C_V}{C_p} \left( \frac{\partial V}{\partial p } \right)_{\tau} }
\end{gathered}
\end{equation}
This is the equation for the \emph{adiabatic compressibility} $\left( \frac{\partial V}{ \partial P} \right)_{\sigma}$ of a body. The inequality $C_p > C_v$ implies that adiabatic compressibility is always smaller in absolute value than isothermal compressibility. Compare this to Eq. (16.14) of Landau and Lifshitz (1980) and Eq. (A-4.10) of Appendix 4, "Summary of Thermodynamic Relations", pp. 577 of Sabersky, et. al. (1998) \cite{SAHG1998}.

The previous equation can be expressed as 
\[
\left(\frac{\partial p}{\partial V }\right)_{\sigma} = \frac{C_p}{C_V} \left( \frac{\partial p}{\partial V } \right)_{\tau} 
\]
From the definition of density $\rho$, $\rho := \frac{MN}{V}$ or $V = \frac{MN}{\rho}$, and assuming $N$ constant, then
\[
\frac{\partial}{\partial V} = \frac{\partial \rho}{\partial V} \frac{\partial }{\partial \rho} = \frac{-MN}{V^2} \frac{\partial}{\partial \rho} = \frac{-\rho^2}{MN} \frac{\partial }{\partial \rho}
\]
so that 
\[
\left(\frac{\partial p}{\partial \rho }\right)_{\sigma} = \frac{C_p}{C_V} \left( \frac{\partial p}{\partial \rho  } \right)_{\tau} 
\]

